\section{The mass kernels of the chiral block}
\label{sec:kernels}

We now return to three constituents and to the block $V_*=\operatorname{im}P_*$ of \cref{sec:pstar}, which carries one Pati--Salam family
$(\mathbf4,\mathbf2,\mathbf1)_{\Gamma=-1}\oplus(\mathbf4,\mathbf1,\mathbf2)_{\Gamma=+1}$ with a right-handed neutrino. A vacuum for this matter is
\emph{dynamically selected} if it solves $\delta\Gamma_{\rm eff}/\delta\Phi=0$ with a second variation that is non-negative apart from the zero modes
of broken symmetries, the equations being those of a stated theory. It is \emph{compatible with matter} if, over it, the fermions carry a
well-defined Dirac operator, the unbroken gauge algebra acts consistently on them, anomalies cancel, and the permitted couplings produce
acceptable masses. This and the next three sections answer both questions inside the arena. Three ingredients are inputs, and we name them
where they enter: the projector $P_*$, which is the field content; the promotion of the space of mass kernels to a dynamical field; and, for
the radiative statements, the form of the tree potential at the matching scale.

\subsection{Notation on the block}

The colour algebra $\so6\cong\su4$ is spanned by the fifteen bivectors of the colour directions $C=\{10,11,12,16,17,18\}$; with
$w=(E_7,E_{13},E_{14},E_{15})$ and $\tau_i=-\tfrac12w_jw_k$ for $(i,j,k)$ a cyclic permutation of $(1,2,3)$, as in \cref{sec:pstar},
the two weak algebras are $t^L_i=\tfrac12(\id-\chi_W)\tau_i$ and
$t^R_i=\tfrac12(\id+\chi_W)\tau_i$. These $15+3+3=21$ generators and the two central elements $J$ and $J\Gamma$ commute with $P_*$, with the
Lorentz action and with $J$, and on the block $q_A=E_2E_3$ acts as $J$ and $E_8E_9$ as $J\Gamma$. With $H=-(E_{10}E_{16}+E_{11}E_{17}+E_{12}E_{18})$
the charges are
\[
\hat t=-Jt,\qquad T^L_3=-Jt^L_3,\quad T^R_3=-Jt^R_3,\qquad \BmL=-\tfrac13JH,\qquad Y=T^R_3+\tfrac12(\BmL),
\]
and the colour $\su3$ is the part of $\so6$ commuting with $H$. The charge table is that of \cref{prop:PSledger}. Finally $\gamma^0$
commutes with all $91$ internal bivectors, with $J$ and with the rotations, and anticommutes with $\Gamma$ and the boosts, so it maps $V_*$
onto an orthogonal rank-$64$ \emph{mirror} on which the same internal algebra acts with $\Gamma$ reversed (\cref{rem:PSsettles}).

\subsection{The kernel space}

\begin{definition}[Mass kernels]
\label{def:kernel}
A Lorentz-invariant Grassmann mass kernel on the block is an $M\in\End(V_*)$ with
\[
M^{\!\top}=-M,\qquad [M,P_*\Sigma^{ab}P_*]=0,\qquad \{M,P_*K^aP_*\}=0 ,
\]
for the interior rotations $\Sigma^{ab}$ and boosts $K^a=\tfrac12\gamma^0\gamma^a$. Their real vector space is $\Kc$, with inner product
$\langle M,N\rangle=\Tr M^{\!\top}N$.
\end{definition}

The definition is forced. For a Grassmann-valued $\psi\in V_*\otimes\Lambda$ the form $\psi^{\!\top}M\psi$ sees only the antisymmetric part of
$M$, and it is invariant under $\delta\psi=X\psi$ iff $X^{\!\top}M+MX=0$; the rotations are antisymmetric and the boosts symmetric, which gives the
two conditions. A kernel is exactly a Lorentz-invariant mass term of the block.

\begin{theorem}[The kernel space]
\label{thm:kernelspace}
\begin{enumerate}[label=(\roman*),leftmargin=2.2em]
\item $\dim_\R\Kc=272$. A basis is given by the compressions $P_*BP_*$ of the blades $B=\gamma^0e_I$ and $B=\gamma^1\gamma^2\gamma^3e_I$, $e_I$ an
internal blade, that are antisymmetric and commute with the three involutions defining $P_*$: there are $2176$ such blades and $272$ distinct,
mutually orthogonal compressions, and they span all solutions of \cref{def:kernel}.
\item Under $\su4\oplus\su2_L\oplus\su2_R$,
\[
\begin{aligned}
\Kc\cong{}&\underbrace{(\mathbf{10},\mathbf3,\mathbf1)}_{60}\oplus\underbrace{(\mathbf{10},\mathbf1,\mathbf3)}_{60}\oplus\underbrace{(\mathbf{15},\mathbf2,\mathbf2)}_{120}\\
&\oplus\underbrace{2(\mathbf6,\mathbf1,\mathbf1)}_{24}\oplus\underbrace{(\mathbf1,\mathbf2,\mathbf2)}_{8},
\end{aligned}
\]
the components being identified by their Casimir labels, $(-36,-2,0)$, $(-36,0,-2)$, $(-32,-\tfrac34,-\tfrac34)$, $(-20,0,0)$ and
$(0,-\tfrac34,-\tfrac34)$ respectively.
\item Every basis kernel either anticommutes with $J$ and commutes with $\Gamma$ (\emph{Majorana type}, pairing fermions of equal chirality, $144$ of
them, spanning $(\mathbf{10},\mathbf3,\mathbf1)\oplus(\mathbf{10},\mathbf1,\mathbf3)\oplus2(\mathbf6,\mathbf1,\mathbf1)$), or commutes with $J$ and anticommutes with
$\Gamma$ (\emph{Dirac type}, pairing opposite chiralities, $128$ of them, spanning $(\mathbf{15},\mathbf2,\mathbf2)\oplus(\mathbf1,\mathbf2,\mathbf2)$).
\item No nonzero kernel is invariant under the $21$ Pati--Salam generators.
\item The kernels invariant under the twelve Standard-Model generators form the two-dimensional space spanned by $\Phi_0$ and $J\Phi_0$, where
\begin{equation}
\label{eq:Phi0}
\Phi_0=P_\nu\,P_*\,\gamma^0E_3E_9E_{13}E_{14}E_{16}E_{17}E_{18}\,P_*\,P_\nu ,
\end{equation}
$P_\nu$ being the projector of the block $V_*$ onto $\Gamma=+1$, $T^R_3=+\tfrac12$, $\BmL=-1$, the right-handed neutrino (on all of $S_3$
the joint eigenspace of these three operators has rank eight). Moreover $\Phi_0^{\!\top}\Phi_0=P_\nu$, of rank
four, $J$ anticommutes with $\Phi_0$, and
\[
[t^R_3,\Phi_0]=J\Phi_0,\qquad [\tfrac13H,\Phi_0]=-2J\Phi_0,\qquad [t_Y,\Phi_0]=0,
\]
so $\Phi_0$ carries $T^R_3=1$, $\BmL=-2$ and $Y=0$. It lies in the $(\mathbf{10},\mathbf1,\mathbf3)$ component.
\end{enumerate}
\end{theorem}

\begin{proof}[Proof of the identification in (iii)]
A kernel commuting with $\Gamma$ preserves the two chirality halves, which carry $(\mathbf4,\mathbf2,\mathbf1)$ and $(\mathbf4,\mathbf1,\mathbf2)$; one anticommuting with
$\Gamma$ exchanges them. A bilinear in two $(\mathbf4,\mathbf2,\mathbf1)$ is an $\su2_R$ singlet and one in two $(\mathbf4,\mathbf1,\mathbf2)$ an $\su2_L$ singlet, while a
bilinear between the halves is a doublet of both. So $(\mathbf{15},\mathbf2,\mathbf2)$ and $(\mathbf1,\mathbf2,\mathbf2)$ are of the second type and the others of
the first, and $60+60+24=144$ and $120+8=128$ agree with the counts.
\end{proof}

\begin{figure}[t]
\centering
\begin{tikzpicture}[x=0.055cm,y=1cm,font=\footnotesize]
\def\kbar#1#2#3#4#5{\fill[#4] (#1,0) rectangle (#2,0.9); \draw (#1,0) rectangle (#2,0.9); \node at ({(#1+#2)/2},0.45) {#3}; \node[below,font=\footnotesize] at ({(#1+#2)/2},0) {#5};}
\kbar{0}{60}{$(\mathbf{10},\mathbf3,\mathbf1)$}{tcol!20}{$60$}
\kbar{60}{120}{$(\mathbf{10},\mathbf1,\mathbf3)$}{tcol!40}{$60$}
\kbar{120}{144}{\scriptsize$2(\mathbf6,\mathbf1,\mathbf1)$}{tcol!10}{$24$}
\kbar{144}{264}{$(\mathbf{15},\mathbf2,\mathbf2)$}{scol!20}{$120$}
\kbar{264}{272}{}{scol!40}{$8$}
\node[anchor=south,font=\footnotesize] at (268,0.95) {$(\mathbf1,\mathbf2,\mathbf2)$};
\draw[decorate,decoration={brace,amplitude=5pt}] (0,1.0) -- node[above=5pt] {Majorana type: anticommute with $J$, $144$} (144,1.0);
\draw[decorate,decoration={brace,amplitude=5pt,mirror}] (144,-0.45) -- node[below=5pt] {Dirac type: anticommute with $\Gamma$, $128$} (272,-0.45);
\fill[gcol] (112,0.16) circle (2.2pt);
\draw[gcol,thick,->] (112,0.16) -- (112,-0.55) node[below,align=center,font=\footnotesize] {$\Phi_0$: the only\\ Standard-Model line};
\end{tikzpicture}
\caption{The $272$-dimensional space of Lorentz-invariant mass kernels of the chiral block (\cref{thm:kernelspace}), drawn to scale. The
Standard-Model invariant kernels form a single complex line, the right-handed neutrino Majorana kernel $\Phi_0$, inside $(\mathbf{10},\mathbf1,\mathbf3)$.}
\label{fig:kernels}
\end{figure}

The complex dimension of $\Kc$ is $136=\dim\mathrm{Sym}^2(\C^{16})$: a Lorentz-invariant mass term of sixteen Weyl fermions is a complex
symmetric $16\times16$ matrix. In the language of $\so{10}$, $\mathrm{Sym}^2(\mathbf{16})=\mathbf{10}\oplus\mathbf{126}$ \cite{Slansky}, and the
decomposition (ii) is the Pati--Salam branching of $\mathbf{10}\oplus\mathbf{126}$; the $(\mathbf{10},\mathbf1,\mathbf3)\subset\mathbf{126}$ is the field which, in
minimal left--right and Pati--Salam models, gives the right-handed neutrino a Majorana mass \cite{MohapatraSenjanovic80,MohapatraSenjanovic81}.
Here this content has been obtained from \cref{def:kernel}, with only the Pati--Salam algebra acting. The only component containing a
Standard-Model singlet is $(\mathbf{10},\mathbf1,\mathbf3)$: the colour-singlet parts of $(\mathbf{10},\mathbf3,\mathbf1)$, $(\mathbf{15},\mathbf2,\mathbf2)$ and $(\mathbf1,\mathbf2,\mathbf2)$ carry
$Y=\pm1$ or $\pm\tfrac12$, and $(\mathbf6,\mathbf1,\mathbf1)$ has none.

\begin{theorem}[Symmetric criticality]
\label{thm:symcrit}
Let $V:\Kc\to\R$ be smooth and invariant under the Pati--Salam group acting by $M\mapsto gMg^{-1}$. Then for every $r\in\R$
\[
\nabla V(r\Phi_0)=\frac1{\langle\Phi_0,\Phi_0\rangle}\,\frac{d}{dr}V(r\Phi_0)\,\Phi_0 ,
\]
so $r\Phi_0$ is a critical point of $V$ if and only if $r$ is a critical point of $r\mapsto V(r\Phi_0)$.
\end{theorem}

\begin{proof}
The action is orthogonal. If $h$ fixes $x=r\Phi_0$, invariance gives $h\nabla V(x)=\nabla V(x)$, so $\nabla V(x)$ is fixed by the Standard-Model
group and lies in $\operatorname{span}\{\Phi_0,J\Phi_0\}$. Since $J\Phi_0=[t^R_3,\Phi_0]$ is tangent to the orbit, $\langle\nabla V(x),J\Phi_0\rangle=0$, and
only the radial component remains.
\end{proof}

This is Michel's theorem on critical orbits \cite{Michel71}, the principle of symmetric criticality \cite{Palais79}. Its content here is that no
invariant coupling on the whole kernel space can generate a tadpole for any other component along $\Phi_0$. It says nothing about the second
variation or about other minima, which are the subject of \cref{sec:selection}.

\subsection{The kernel field and its action}

\begin{assumption}
\label{ass:field}
The kernels are the values of a scalar field $\Phi(x)\in\Kc$ with kinetic term $\tfrac12\langle D_\mu\Phi,D^\mu\Phi\rangle$ and a potential $V(\Phi)$
invariant under the gauge algebra, and a single Yukawa coupling $y$ in \eqref{eq:Saction}. The overall normalisation of the kinetic term is a
convention; its use of one trace form on all of $\Kc$, the single coupling $y$ and the form of $V$ are choices, restricted below only by symmetry.
\end{assumption}

With a real Grassmann field $\psi(x)\in V_*$, a gauge field $A_\mu$ with values in an algebra $\mathfrak g$ of antisymmetric operators, the kinetic
kernels $K^\mu=P_*\gamma^0\gamma^\mu P_*$ (symmetric, $K^0=\id$ on $V_*$) and the module metric $G_{ab}=\Tr(T_a^{\!\top}T_b)/512$, take
\begin{equation}
\label{eq:Saction}
S=\int d^4x\Bigl[\tfrac12\psi^{\!\top}K^\mu D_\mu\psi+\tfrac y2\psi^{\!\top}\Phi\psi+\tfrac12\langle D_\mu\Phi,D^\mu\Phi\rangle-V(\Phi)-\tfrac14G_{ab}F^a_{\mu\nu}F^{b\mu\nu}\Bigr],
\end{equation}
with $D_\mu\psi=\partial_\mu\psi+A_\mu\psi$ and $D_\mu\Phi=\partial_\mu\Phi+[A_\mu,\Phi]$. Every term is Lorentz invariant by \cref{def:kernel}, and gauge
invariant for any $\mathfrak g$ inside the $23$-dimensional faithful algebra of the block. The module metric is $1$ on the colour bivectors and
$\tfrac18$ on $t^L$ and $t^R$.

\begin{proposition}[Fermion number]
\label{prop:fermionnumber}
$J$ commutes with the kinetic kernels and with the $21$ Pati--Salam generators. The transformation $\psi\mapsto e^{\theta J}\psi$,
$\Phi\mapsto e^{\theta J}\Phi e^{-\theta J}$, $A\mapsto A$ leaves every term of \eqref{eq:Saction} except $V$ invariant; on $\Kc$ it acts as
$\Phi\mapsto e^{2\theta J}\Phi$ on Majorana-type kernels and trivially on Dirac-type ones. In the complex structure $J$ each of the sixteen Weyl
fermions has charge $1$, so this is the fermion number $F=3B+L$, and $4B=F+(\BmL)$.
\end{proposition}

By \cref{prop:PSanom} the two central directions $J$ and $J\Gamma$ are anomalous on this matter, so the gauge algebra can be any anomaly-free
subalgebra of the $23$, and $\su4\oplus\su2_L\oplus\su2_R$ is the largest one containing Pati--Salam. A gauge field coupled to an anomalous current
makes the quantum theory inconsistent \cite{GrossJackiw72,BouchiatIliopoulosMeyer72}, which is why the central directions carry no gauge
field; they remain global symmetries. In the vacuum $r\Phi_0$ the stabiliser among the $23$ generators is fourteen-dimensional: the twelve
Standard-Model generators together with $J+\tfrac13H$ and $J-J\Gamma$, whose charges are $1+(\BmL)=4B$ and $1-\Gamma$, baryon number and twice
the left-handed fermion number, both global.

\begin{proposition}[The coupled field equations]
\label{prop:coupled3}
The Euler--Lagrange equations of \eqref{eq:Saction} are
\[
K^\mu D_\mu\psi+y\Phi\psi=0,\qquad D_\mu D^\mu\Phi+\nabla V(\Phi)-\tfrac y2\Pi_\Kc(\psi\psi^{\!\top})=0,\qquad G_{ab}D_\nu F^{b\nu\mu}=-\bigl(j^\mu_a(\psi)+j^\mu_a(\Phi)\bigr),
\]
with $\Pi_\Kc$ the orthogonal projection onto $\Kc$. At constant $\Phi=r\Phi_0$, $\psi=0$, $A=0$ the first and third hold identically and the second
reduces to $\tfrac{d}{dr}V(r\Phi_0)=0$ by \cref{thm:symcrit}. The same holds with $V$ replaced by the effective potential, which is again
Pati--Salam invariant; at one loop
\begin{equation}
\label{eq:V1}
V_1=\frac1{64\pi^2}\,\mathrm{STr}\Bigl[\mathcal M^4\Bigl(\ln\frac{\mathcal M^2}{\mu_R^2}-c\Bigr)\Bigr],\qquad
\mathrm{STr}=\textstyle\sum_{\rm scalars}+3\sum_{\rm vectors}-2\sum_{\rm Weyl},
\end{equation}
the fermion term being the determinant of the Dirac operator of the first equation \cite{ColemanWeinberg73,Jackiw74}; at stationary points it
does not depend on the gauge fixing \cite{Nielsen75}.
\end{proposition}

The fermions, classically zero, act on the vacuum only through the effective potential, and the matter content enters through the Yukawa
coupling $y$. That there is one such coupling is part of \cref{ass:field}, not a consequence of calling $\Phi$ one field: $\Phi$ is reducible under the
Pati--Salam algebra, and symmetry alone allows an independent Yukawa coupling and an independent kinetic normalisation for each of its five irreducible
components, with a mixing between the two copies of $(\mathbf6,\mathbf1,\mathbf1)$. The same is true of the common gauge coupling of $\su4$ and $\su2_R$ used in
\cref{sec:selection}. In the vacuum the block Hamiltonian is $H(p)=p_iK^i+iy\,r\,\Phi_0$; since $\{K^i,K^j\}=2\delta^{ij}$,
$\{K^i,\Phi_0\}=0$ and $-\Phi_0^2=P_\nu$,
\[
H(p)^2=|p|^2+y^2r^2P_\nu :
\]
one Majorana fermion, the right-handed neutrino, with four real components and mass $|y|r$, and sixty massless real components.

\Cref{thm:symcrit} settles stationarity for the whole coupled quantum system along $\Phi_0$. Which orbit is the \emph{minimum} is decided by the
couplings, and that is the subject of the next section. A non-negative second variation on the $(\mathbf{10},\mathbf1,\mathbf3)$ directions alone does not
imply it on all of $\Kc$: fluctuations outside $(\mathbf{10},\mathbf1,\mathbf3)$ with the Standard-Model quantum numbers of $(\mathbf{10},\mathbf1,\mathbf3)$ fluctuations can
mix with them, and \cref{sec:fullK} exhibits such a mixing.
